% Policy Indistinguishability section body.
\subsection{Policy Indistinguishability (PI)}
\label{sec:eval-pi}

\paragraph{Property.} PI strengthens PCH: even when two policies are \emph{semantically
related}, the observer must not tell which one was executed---formally,
$\Pr[\hat b = b]\approx \tfrac12$. We test four pair archetypes: \texttt{semantic\_%
equivalent} (identical normalized logic), \texttt{structural\_variant} (same semantics,
different execution protections), \texttt{weak\_equivalent} ($\ge 0.95$ semantic
overlap), and \texttt{controlled\_difference} (small targeted deltas; a regression
control).

\paragraph{Attacker levels.} The distinguisher's task is to decide, from the on-chain
features of an execution, which policy of the pair produced it. We try three attackers of
increasing strength: \emph{Level~I} trains a classifier on the features of a single
execution; \emph{Level~II} compares the feature \emph{difference} between the two
policies' executions (a paired attacker); and \emph{Level~III} runs two-sample
statistical tests (Kolmogorov--Smirnov and Maximum Mean Discrepancy) directly on the two
feature distributions. A pair is indistinguishable only if all three stay at chance.

\paragraph{Results.} Our primary attacker is Level~I, evaluated over 30 policy-disjoint
seeds (Table~\ref{tab:pi-pairs}): the overall AUC is $0.480\pm0.042$ with $0\%$ of seeds
significant, and every archetype that has enough test data passes. Only
\texttt{weak\_equivalent} is set aside---its median test set is $n=9$, below the
power threshold of Section~\ref{sec:eval-method}, so it is reported but not gated. The
stronger Level~II and Level~III attackers, measured on a single split as a diagnostic,
are likewise at chance (Level~II overall AUC $0.500$; Level~III MMD $\le 0.05$ with
non-significant KS), so no attacker separates the pairs.

\paragraph{Why multi-seed matters.} At three seeds the \texttt{structural\_variant}
archetype read AUC $0.69$ with $100\%$ of seeds significant---enough to look like a
leak---but at 30 seeds it regressed to $0.541\pm0.095$ with only $40\%$ significant,
revealing the spike as a small-sample fluke. The final 30-seed verdict
(Table~\ref{tab:pi-pairs}, AUC $0.535\pm0.088$, $20\%$ significant) confirms the pass.
This is precisely the failure mode the multi-seed criterion of
Section~\ref{sec:eval-method} is designed to reject.

\paragraph{Clustering attack.} On balanced $30$-vs-$30$ pairs (133 pairs, after the
sample-size correction of Section~\ref{sec:eval-method}) K-Means recovers no structure:
mean ARI $0.0015$, mean observed NMI $0.0136$, mean debiased NMI $0.0012$
(Table~\ref{tab:clustering-summary}). Only $6/133$ pairs cross the per-pair significance
threshold---consistent with the $\approx 5\%$ false-positive rate expected by chance at
$\alpha{=}0.05$---so PI holds under the unsupervised attacker as well.
